\documentclass{exam}
\usepackage{math-macros}

\department{}
\school{}
\teacher{Vũ Vương An}
\topic{Ứng dụng đạo hàm để khảo sát và vẽ đồ thị hàm số}
\examtitle{ĐÁP ÁN — PHIẾU BÀI TẬP}
\subject{Toán}
\grade{12}
\duration{20 phút}
\examdate{17/09/2026}
\setquestiontotal{ 10 }

\answerkeytrue


\begin{document}
\makeheader


\begin{question}
Hàm số $y = x^4 - 2x^2$ có bao nhiêu điểm cực tiểu?



\choicegrid{ $0$ }{ $1$ }{ $2$ }{ $3$ }


\end{question}

\begin{dapan}
$2$
\end{dapan}


\begin{loigiai}
$y' = 4x^3 - 4x = 4x(x-1)(x+1)$. Các điểm tới hạn $x = -1,\,0,\,1$. $y'' = 12x^2 - 4$: $y''(\pm 1) > 0$ (cực tiểu), $y''(0) < 0$ (cực đại). Vậy có hai điểm cực tiểu.
\end{loigiai}



\begin{question}
Cho hàm số $y = x^3 - 3x^2 + 2$. Hoành độ điểm cực đại của đồ thị hàm số là



\choicegrid{ $-1$ }{ $0$ }{ $1$ }{ $2$ }


\end{question}

\begin{dapan}
$0$
\end{dapan}


\begin{loigiai}
Ta có $y' = 3x^2 - 6x = 3x(x-2)$. Các điểm tới hạn $x = 0$ và $x = 2$. Bảng biến thiên: $y'$ đổi dấu từ $+$ sang $-$ tại $x = 0$ nên đồ thị đạt cực đại tại $x = 0$.
\end{loigiai}



\begin{question}
Hàm số $y = -x^3 + 3x$ có bao nhiêu điểm cực trị?



\choicegrid{ $0$ }{ $1$ }{ $2$ }{ $3$ }


\end{question}

\begin{dapan}
$2$
\end{dapan}


\begin{loigiai}
$y' = -3x^2 + 3 = 3(1-x^2)$. $y' = 0$ khi $x = \pm 1$. Vậy hàm số có hai điểm cực trị.
\end{loigiai}



\begin{question}
Hàm số $y = x^3 + 3x^2$ nghịch biến trên khoảng nào dưới đây?



\choicegrid{ $(-\infty; -2)$ }{ $(-2; 0)$ }{ $(0; +\infty)$ }{ $(-1; 1)$ }


\end{question}

\begin{dapan}
$(-2; 0)$
\end{dapan}


\begin{loigiai}
Ta có $y' = 3x^2 + 6x = 3x(x+2)$. $y' < 0$ khi $-2 < x < 0$, nên hàm số nghịch biến trên $(-2; 0)$.
\end{loigiai}



\begin{question}
Giá trị nhỏ nhất của hàm số $y = x^3 - 3x + 2$ trên đoạn $[0; 2]$ bằng



\choicegrid{ $-2$ }{ $-1$ }{ $0$ }{ $1$ }


\end{question}

\begin{dapan}
$0$
\end{dapan}


\begin{loigiai}
$y' = 3x^2 - 3 = 3(x-1)(x+1)$. Trên $[0; 2]$ xét $x = 0,\,1,\,2$: $y(0) = 2$, $y(1) = 0$, $y(2) = 4$. Vậy $\min_{[0;2]} y = 0$.
\end{loigiai}



\begin{question}
Giá trị lớn nhất của hàm số $y = x^3 - 3x + 1$ trên đoạn $[-1; 2]$ bằng



\choicegrid{ $1$ }{ $2$ }{ $3$ }{ $5$ }


\end{question}

\begin{dapan}
$3$
\end{dapan}


\begin{loigiai}
$y' = 3x^2 - 3 = 3(x-1)(x+1)$. Xét $x \in [-1; 2]$: $y(-1) = 3$, $y(1) = -1$, $y(2) = 3$. Vậy $\max_{[-1;2]} y = 3$.
\end{loigiai}



\begin{question}
Hàm số $y = x^3 - 3x$ đồng biến trên khoảng nào dưới đây?



\choicegrid{ $(-1; 1)$ }{ $(0; 2)$ }{ $(1; +\infty)$ }{ $(-2; 0)$ }


\end{question}

\begin{dapan}
$(1; +\infty)$
\end{dapan}


\begin{loigiai}
$y' = 3x^2 - 3 = 3(x-1)(x+1)$. $y' > 0$ khi $x < -1$ hoặc $x > 1$. Trong các phương án, khoảng $(1; +\infty)$ là khoảng đồng biến.
\end{loigiai}



\begin{question}
Một chất điểm chuyển động trên trục $Ox$ với vận tốc $v(t) = 3t^2 - 12t + 9$ (m/s), $t$ tính bằng giây. Thời điểm chất điểm đứng yên lần đầu tiên là



\choicegrid{ $1$ s }{ $2$ s }{ $3$ s }{ $4$ s }


\end{question}

\begin{dapan}
$1$ s
\end{dapan}


\begin{loigiai}
Chất điểm đứng yên khi $v(t) = 0 \Leftrightarrow 3(t^2 - 4t + 3) = 0 \Leftrightarrow (t-1)(t-3) = 0$. Các nghiệm $t = 1$ và $t = 3$. Thời điểm đầu tiên là $t = 1$ s.
\end{loigiai}



\begin{question}
Tìm $m$ để hàm số $y = x^3 - 3x^2 + m - 1$ đạt cực đại bằng $0$.



\choicegrid{ $0$ }{ $1$ }{ $2$ }{ $3$ }


\end{question}

\begin{dapan}
$1$
\end{dapan}


\begin{loigiai}
$y' = 3x^2 - 6x = 3x(x-2)$. Hàm số đạt cực đại tại $x = 0$, $y(0) = m - 1$. Yêu cầu $m - 1 = 0 \Leftrightarrow m = 1$.
\end{loigiai}



\begin{question}
Đồ thị hàm số $y = x^3 - 3x$ cắt đường thẳng $y = m$ tại ba điểm phân biệt khi



\choicegrid{ $m = 2$ }{ $m = -2$ }{ $m = 0$ }{ $m = 3$ }


\end{question}

\begin{dapan}
$m = 0$
\end{dapan}


\begin{loigiai}
Đặt $f(x) = x^3 - 3x$. Cực đại $f(-1) = 2$, cực tiểu $f(1) = -2$. Đường thẳng $y = m$ cắt đồ thị tại ba điểm phân biệt khi $-2 < m < 2$. Trong các phương án chỉ có $m = 0$ thỏa.
\end{loigiai}




\end{document}
